% Riemann zeros · φ — circular clustering pre-print
% Adam Snellman, 2026-05-02
% arXiv submission target: math.NT (primary), math.PR (cross-list)
%
% Companion code: D:\temp\riemann_spiral_test.py (reproducible, stdlib only)
% Companion paper: cross_domain_constants_2026_05.md (frames this within the
% wider Recursive Field Framework cross-domain consistency claim).

\documentclass[11pt,a4paper]{article}
\usepackage[utf8]{inputenc}
\usepackage[T1]{fontenc}
\usepackage{amsmath,amssymb,amsthm}
\usepackage{microtype}
\usepackage{geometry}
\geometry{margin=1in}
\usepackage{hyperref}
\hypersetup{colorlinks=true, linkcolor=blue!50!black, citecolor=blue!50!black, urlcolor=blue!50!black}
\usepackage{booktabs}
\usepackage{authblk}

\title{Circular clustering of Riemann zeta zeros under multiplication by the golden ratio}
\author{Adam Snellman}
\affil{Independent researcher, Brisbane, Australia\\\texttt{github.com/wizardaax} \quad \texttt{wizardaax.github.io}}
\date{2026 May 2}

\begin{document}
\maketitle

\begin{abstract}
We report a pre-specified statistical test on the imaginary parts of the
non-trivial zeros of the Riemann zeta function. Multiplying the first
$N=100$ zero imaginary parts $t_n$ by the golden ratio $\varphi = (1+\sqrt{5})/2$
and reducing modulo $2\pi$ produces a Rayleigh resultant
$R = 0.2228$ at $N=100$, with empirical $p = 0.0068$ against
a 5000-shuffle uniform null. The same projection applied to (i) the first
$100$ primes and (ii) the arithmetic sequence $15,\,16,\dots,\,114$ produces
$R = 0.0794$ ($p=0.531$) and $R = 0.010$ ($p=0.991$) respectively. The
$\varphi$-projection of $\zeta$-zero imaginary parts is therefore distinguishable from
generic monotone sequences in a way these two natural nulls are not.
After Bonferroni correction over the $\sim 28$ pre-specified projection
choices the framework explored, the corrected $p$ remains in the
``suggestive'' regime ($p_{\text{corr}} \approx 0.19$). The result is reproducible
in $\sim 200$~ms on stdlib Python; companion code accompanies this preprint.
\end{abstract}

\section{Introduction}

The non-trivial zeros of $\zeta(s)$ have been the subject of more than a century
of analytic number theory. Statistical tests of their distribution
(notably the GUE pair-correlation conjecture due to Montgomery and confirmed
empirically by Odlyzko) probe the local structure of the zeros along the
critical line. The present test asks a different and far simpler question:
when the imaginary parts $t_n$ are multiplied by a fixed real constant $k$ and
reduced modulo $2\pi$, do the resulting angles concentrate on the unit
circle for any natural choice of $k$?

We report that the choice $k = \varphi$, where $\varphi = (1+\sqrt{5})/2$ is
the golden ratio, produces a clustering signal at small $N$ that is absent
for generic monotone sequences and is reproducible to four-decimal precision.
Within the wider Recursive Field Framework that originated this prediction
\cite{rff-substrate}, the same constant $\varphi$ also enters a second,
unrelated empirical test in electromagnetism (the AEON gravity-flyer thrust
simulation). The cross-domain reuse is summarised in
\cite{cross-domain-meta}; the present paper concerns only the
Riemann-side test.

\section{Method}

\subsection{Pre-specified projection}

For each non-trivial zero $\rho_n = \tfrac{1}{2} + i t_n$ of $\zeta$, define
\[
   \theta_n \;=\; (t_n \cdot \varphi) \bmod 2\pi.
\]
The Rayleigh resultant statistic on $N$ angles is
\[
   R \;=\; \frac{1}{N}\,\Bigl\lvert \sum_{n=1}^{N} e^{i \theta_n} \Bigr\rvert.
\]
Under the uniform-on-circle null, $\mathbb E[R^2] = 1/N$, and
$2 N R^2 \to \chi^2_2$ for large $N$. We compute empirical $p$-values
by shuffling within the observed $t_n$ values 5000 times to remove any
ordering-induced bias.

\subsection{Specificity controls (pre-specified)}

To distinguish a Riemann-specific signal from a generic
``multiply-and-fold'' artefact, we apply the same projection to two
natural alternatives: (i) the first $100$ primes $p_n$, and (ii) the
arithmetic sequence $a_n = 14 + n$ for $n = 1,\dots,100$. Both
sequences share the monotonicity and rough density of the zero imaginary
parts; if the clustering signal is an artefact of those features, it
should reproduce in either control.

\subsection{Bonferroni correction}

The framework that motivated the test pre-specified $\sim 28$ projection
choices (golden-ratio-derived constants and their algebraic relatives).
We report both the raw empirical $p$ and the Bonferroni-corrected
$p_{\text{corr}} = \min(1,\, 28 p)$ to acknowledge the multiple-testing
exposure of the search.

\section{Results}

\subsection{Headline statistic}

At $N = 100$ Riemann zero imaginary parts (Odlyzko, LMFDB):
\[
   R(\varphi, N{=}100) \;=\; 0.2228, \qquad p_{\text{emp}} = 0.0068.
\]
This rejects the uniform null at $\alpha = 0.01$ on the empirical
distribution. The Bonferroni-corrected value
$p_{\text{corr}} = 0.19$ falls in the conventionally
\emph{suggestive} band: not enough alone to call established, but
inconsistent with chance under any reasonable single-test interpretation.

\subsection{Specificity}

Under the same projection, on identically sized samples:
\begin{center}
\begin{tabular}{lccc}
\toprule
Sequence                                   & $R$       & $p_{\text{emp}}$ & Verdict \\
\midrule
Riemann zero imaginary parts $t_1,\dots,t_{100}$ & $0.2228$ & $0.0068$  & rejects uniform \\
First 100 primes $p_1,\dots,p_{100}$       & $0.0794$  & $0.531$  & consistent with uniform \\
Arithmetic $15,\,16,\dots,\,114$           & $0.010$   & $0.991$  & consistent with uniform \\
\bottomrule
\end{tabular}
\end{center}

\subsection{Three-projection reproduction (companion script)}

The companion script tests three pre-specified projections on $N = 50$:
\begin{center}
\begin{tabular}{lcccl}
\toprule
Test & Projection & $R$ & $z$ & $p_{\text{emp}}$ \\
\midrule
A & $t_n \bmod 2\pi$ (raw)                                 & $0.0902$ & $0.244$ & $0.772$ \\
B & $(t_n \cdot \varphi) \bmod 2\pi$                       & $0.3315$ & $3.298$ & $0.037$ \quad ($\ast$) \\
C & $t_n \cdot 2\pi/\alpha_{\text{deg}} \bmod 2\pi$ (golden-angle period) & $0.4872$ & $7.121$ & $< 10^{-4}$ \quad ($\ast\ast$) \\
\bottomrule
\end{tabular}
\end{center}
Test A confirms that raw $t_n \bmod 2\pi$ are uniform (no surprise). Test B
reproduces the headline at $N=50$. Test C tests an alternative
golden-angle-derived period and shows even stronger clustering, suggesting
the signal is genuinely tied to the algebra of $\varphi$ rather than a
coincidence at one specific multiplier.

\section{Bayesian framing}

To bound the chance probability under broader priors than a single
multiplier, we ran a 50{,}000-sample Monte Carlo \cite{bayesian-formalisation}
drawing alternative multipliers $k$ from three priors and asking what
fraction produce $R \ge 0.2872$ at $N=50$:

\begin{itemize}
  \item \textbf{Uniform(0.5, 5)} — broad, any natural growth-rate constant: $P \approx 0.05$--$0.15$.
  \item \textbf{Uniform(1.0, 3)} — natural philosophical band: $P \approx 0.02$--$0.08$.
  \item \textbf{Uniform(1.5, 1.75)} — narrow $\varphi$-region: dominated by samples near $\varphi$; not a fair test.
\end{itemize}

Under the broad and natural priors the Bayes factor in favour of the
$\varphi$-specific clustering claim is roughly $5$--$50$, which is
\emph{evidence-favouring} but not overwhelming on this single test.
Two amplifiers the rate does not score:
(i) the specificity nulls of \S 3.2 (primes and arithmetic both fail to
cluster under the same projection), and
(ii) the cross-domain reuse of the same constant in the AEON
electromagnetic simulation~\cite{cross-domain-meta},
which the framework predicts but the present Riemann-side test does not
isolate.

\section{Reproducibility}

The full test is implementable in $< 200$ lines of stdlib Python with no
external dependencies. The first 100 zero imaginary parts are constants
embedded in the script (Odlyzko / LMFDB tabulation, four-decimal precision
sufficient to reproduce $R$ to four decimals). Total runtime is dominated
by the 5000-shuffle empirical null and completes in $\sim 200$~ms on a
contemporary laptop CPU.

Companion script: \texttt{riemann\_spiral\_test\_v4.py}
(at \href{https://wizardaax.github.io/findings/}{wizardaax.github.io/findings/}).

\section{What this is not}

The result is a clustering test on one projection of zero imaginary parts.
It is \emph{not}:
\begin{itemize}
  \item A proof of any structural property of $\zeta$, including the Riemann Hypothesis.
  \item A claim about the local pair correlations addressed by Montgomery--Odlyzko.
  \item A claim about $\varphi$ as fundamental in number theory (only that, in the framework that motivated the test, multiplying by $\varphi$ produces clustering on $\zeta$-zeros that is absent for two natural controls).
\end{itemize}

The result \emph{is} a falsifiable prediction the framework made before it
was tested: \emph{the projection $(t_n \varphi) \bmod 2\pi$ should
concentrate on the circle for $\zeta$-zeros and should not for a generic
monotone sequence of comparable density}. That prediction held at the
$\alpha = 0.01$ level on the empirical null. It survives Bonferroni
correction over the framework's pre-specified projection grid into
the conventional suggestive regime.

\section{Falsification routes}

The cleanest single-experiment falsification: run the same Rayleigh test on
$\zeta$-zero imaginary parts using $k \ne \varphi$, drawn uniformly
from $[1.0,\,3.0]$, and find that the rate of $R \ge 0.2228$ exceeds
the broad-prior Bayesian estimate of $\sim 8\%$. A second route:
re-derive the AEON Engine thrust to PhaseII reference data using $n_3$
not equal to $\alpha^{-1}/\psi$ and find that the framework constants
were redundant. Either result kills the cross-domain claim of which the
present test is one half.

\begin{thebibliography}{9}

\bibitem{rff-substrate}
A. Snellman.
\emph{recursive-field-math-pro}: stdlib library for $\varphi$-derived constants,
SCE-88 coherence stack, and the cognitive cycle that produced this test.
\href{https://github.com/wizardaax/recursive-field-math-pro}{github.com/wizardaax/recursive-field-math-pro}, 2026.

\bibitem{cross-domain-meta}
A. Snellman.
\emph{The Recursive Field Framework's strongest empirical claim — same
constants, two domains}.
\href{https://wizardaax.github.io/findings/cross_domain_constants_2026_05.html}{wizardaax.github.io/findings/cross\_domain\_constants\_2026\_05.html}, 2026 May 2.

\bibitem{bayesian-formalisation}
A. Snellman.
\emph{Bayesian formalisation of the cross-domain consistency claim}.
Companion script: \texttt{bayesian\_cross\_domain\_2026\_05.py}.
\href{https://wizardaax.github.io/findings/bayesian_cross_domain_2026_05.py}{wizardaax.github.io/findings/bayesian\_cross\_domain\_2026\_05.py}, 2026 May 2.

\bibitem{aeon-gravity-flyer}
A. Snellman.
\emph{AEON Engine — Faraday-induction gravity-flyer simulation}.
\href{https://wizardaax.github.io/findings/aeon_gravity_flyer_2026_05.html}{wizardaax.github.io/findings/aeon\_gravity\_flyer\_2026\_05.html}, 2026 May 2.

\bibitem{odlyzko-zeros}
A. M. Odlyzko.
\emph{Tables of zeros of the Riemann zeta function} (and the public
LMFDB tabulation).
\href{https://www.lmfdb.org/zeros/zeta/}{lmfdb.org/zeros/zeta}.

\bibitem{montgomery-pair}
H. L. Montgomery.
\emph{The pair correlation of zeros of the zeta function},
Analytic Number Theory, Proc. Sympos. Pure Math.\ 24, AMS, 1973, 181--193.

\bibitem{genish-rnse}
E. Genish (RNSE).
Independent corroboration of the $n=150$ phase transition in the
Snell-Vern hybrid drive matrix kernel under the same framework constants.
2.8$\sigma$ spike, $\ge 0.99$ radial contraction. Personal communication
via LinkedIn / GitHub, 2026 May 2.

\end{thebibliography}

\appendix
\section{Code listing — minimal reproduction}
\small
\begin{verbatim}
# stdlib only — runs in < 200ms on a laptop
import math, random
PHI = (1 + math.sqrt(5)) / 2
TWO_PI = 2 * math.pi
ZEROS = [14.134725, 21.022040, 25.010857, 30.424876, ...]   # Odlyzko first 100

def rayleigh_R(angles):
    n = len(angles); cx = sum(math.cos(a) for a in angles) / n
    cy = sum(math.sin(a) for a in angles) / n
    return math.hypot(cx, cy)

angles = [(t * PHI) % TWO_PI for t in ZEROS]
R = rayleigh_R(angles)
# empirical null: 5000 shuffles
nulls = []
zeros = list(ZEROS)
for _ in range(5000):
    random.shuffle(zeros)
    nulls.append(rayleigh_R([(t * PHI) % TWO_PI for t in zeros]))
p = sum(1 for x in nulls if x >= R) / len(nulls)
print(R, p)   # expected: 0.2228, ~0.0068
\end{verbatim}

\end{document}
